% ═══════════════════════════════════════════════════════════════════════════════
%  CAPÍTULO — Factory Method
% ═══════════════════════════════════════════════════════════════════════════════

\chapter{Factory Method}

\patroninfo{Creacional}{107}

% ─── Situación Problema ──────────────────────────────────────────────────────
\section{Situación Problema}

Airbnb necesita enviar notificaciones a sus usuarios ante distintos
eventos: una reserva confirmada, un mensaje de un anfitrión, un recordatorio
de check-in o una alerta de pago pendiente. Dependiendo de las preferencias
del usuario y del tipo de evento, la notificación puede enviarse por correo
electrónico, mensaje push en la aplicación móvil o SMS.

El módulo de eventos actualmente contiene bloques \texttt{if-else} que
deciden qué objeto de notificación instanciar. Cada vez que se incorpora
un nuevo canal —por ejemplo, notificaciones en WhatsApp— es necesario
modificar ese módulo central, rompiéndose el principio de abierto/cerrado
y aumentando el riesgo de introducir errores en la lógica ya probada.
Además, el módulo de eventos no debería tener que conocer los detalles
de construcción de cada tipo de notificación.

% ─── Solución ────────────────────────────────────────────────────────────────
\section{Solución}

El patrón Factory Method delega la responsabilidad de crear el objeto
de notificación a subclases especializadas. Se define una clase base
\texttt{NotificadorEvento} con un método fábrica abstracto que devuelve
un objeto de tipo \texttt{Notificacion}; cada subclase concreta
(\texttt{NotificadorEmail}, \texttt{NotificadorPush}, \texttt{NotificadorSMS})
implementa ese método y determina qué objeto instanciar.

El módulo de eventos trabaja exclusivamente con la interfaz abstracta y
nunca sabe qué tipo concreto está usando. Cuando se requiere un nuevo
canal de comunicación, basta con añadir una nueva subclase sin tocar el
código existente. Esto hace que el sistema sea extensible, más fácil de
probar y libre de la proliferación de condicionales en el núcleo del negocio.

% ─── Diagrama de clases ─────────────────────────────────────────────────────
\begin{figure}[H]
    \centering
    % \includegraphics[width=\textwidth]{imagenes/abstract_factory_diagrama.png}
    \caption{Diagrama de clases del patrón Abstract Factory}
\end{figure}


% ─── Roles de las Clases ────────────────────────────────────────────────────
\section{Roles de las Clases}

% Explicar la responsabilidad de cada clase/interfaz

\begin{itemize}
    \item \textbf{AbstractFactory:}
    \item \textbf{ConcreteFactory:}
    \item \textbf{AbstractProduct:}
    \item \textbf{ConcreteProduct:}
    \item \textbf{Client:}
\end{itemize}


% ─── Main ───────────────────────────────────────────────────────────────────
\section{Main}

% Explicar qué realiza el programa principal

\begin{lstlisting}[language=Java, caption=NombreClase]
 
\end{lstlisting}


% ─── Salida del Main ────────────────────────────────────────────────────────
\section{Salida del Main}

\begin{lstlisting}[language=Java, caption=NombreClase]
 
\end{lstlisting}


% ─── Clases ─────────────────────────────────────────────────────────────────
\section{Clases}

% ─── Clase 1 ────────────────────────────────────────────────────────────────
\subsection{NombreClase}

\begin{lstlisting}[language=Java, caption=NombreClase]
 
\end{lstlisting}


% ─── Clase 2 ────────────────────────────────────────────────────────────────
\subsection{NombreClase}

\begin{lstlisting}[language=Java, caption=NombreClase]
 
\end{lstlisting}


% ─── Clase 3 ────────────────────────────────────────────────────────────────
\subsection{NombreClase}

\begin{lstlisting}[language=Java, caption=NombreClase]
 
\end{lstlisting}


% ─── Conclusiones ───────────────────────────────────────────────────────────
\section{Conclusiones}

% Explicar:
% - Ventajas del patrón
% - Cuándo usarlo
% - Beneficios obtenidos
% - Posibles desventajas
% - Relación con principios SOLID